\begingroup
\sloppy
\section{Unified Mobile Device Mesh, Television and IoT Control}
\label{sec:pilot-unified-mobile-device-mesh}

The unified Pilot mobile application now projects the existing AION Device Fabric into two
purpose-specific Personal surfaces: Devices for mesh visibility and bounded IoT actions, and TV
for shared-screen ownership and verified television control.  The implementation reuses the
Fabric node ledger, discovery engine, signed capability capsules, webOS gateway, television-room
registry, Production Private Identity and local infrared climate gateway.  It does not create a
second device registry or bypass the device adapters already governed by the mother brain.

\subsection{Private Mesh Projection}

The Devices surface shows each evidenced node's name, class, platform, role, enrollment state,
transports and count of signed capabilities.  The mother strips public keys, pairing secrets,
local IP addresses, control endpoints, schema internals and room endpoints before returning the
projection.  A room contains only its display name, television name, node reference and default
marker.  The phone therefore receives enough information to understand the home without becoming
a copy of its network or trust stores.

An explicit ``Discover nearby devices'' action runs the existing bounded discovery engine.  It
listens for standard advertisements, performs the already governed bounded discovery probes and
refreshes schema evidence.  The mobile result declares that enrollment, control and credential
attempts remain false.  A discovered node is visibly labelled observation-only.  Discovery never
creates a capability lease and the phone cannot enroll a new device from this surface.

\subsection{Exclusive Shared-Screen Session}

The television is a shared household surface, while private files, tasks, accounts and business
work belong to a person.  Before any mutating TV command, the person must acquire a five-minute
idle session from a trusted phone.  The browser signs a dedicated possession payload with its
non-exportable Ed25519 key and then signs the enclosing mobile request.  The mother checks both
proofs against the same certificate, person and device.

Only one person may hold the screen.  Another household member sees that the television is in use
but not the controller's private identity.  Every accepted command renews the idle deadline.
Expiry or explicit ``Lock TV workspace'' immediately removes control and hides private shared-
screen content.  An observe-state operation remains available without acquisition because it
changes no television state.

\subsection{Governed Television Controls}

The TV surface provides directional navigation, OK, Back, Home, volume, mute, play and pause, plus
explicit launch controls for Pilot Home, Netflix, YouTube, Games and God View.  The phone cannot
submit arbitrary webOS methods, application identifiers, URLs or pointer commands through this
interface.  Each request is constrained to an allowlist and routed through the existing signed
webOS capability and adapter.

Mutating commands require \texttt{tv.control}; live observation requires the narrower
\texttt{tv.observe}.  A content-bound request identifier is passed into the verified-navigation
ledger so a network retry does not repeat an already accepted command.  The response distinguishes
adapter acceptance, a device receipt and verified navigation.  Button delivery alone is never
reported as successful screen navigation.  The user can request a fresh webOS foreground-app and
audio-state observation, while visually dependent navigation still requires the existing Observer
workflow to prove the screen outcome.

\subsection{IoT Capability Boundary}

The initial non-TV control surface exposes only presets already learned by the local infrared
climate gateway.  No raw infrared waveform, device MAC address or gateway host is sent to the
phone.  Sending a preset requires a signed \texttt{devices.control} capability and an explicit
confirmation.  The response separates \texttt{transport\_delivered} from
\texttt{device\_state\_verified}.  Since ordinary infrared is one-way, Pilot always leaves the
latter false and instructs the person to confirm the appliance or add an authorized state sensor.

Other discovered appliances remain non-controllable until they are enrolled and receive an
appropriate signed capability capsule and adapter.  A generic device name or inferred schema is
not authority.

\subsection{Mobile Capability Scopes}

The production phone lease now separates:
\begin{itemize}
  \item \texttt{devices.read} for the sanitized mesh and room projection;
  \item \texttt{devices.discover} for a bounded user-requested local scan;
  \item \texttt{devices.control} for explicitly supported enrolled IoT actions;
  \item \texttt{tv.observe} for non-mutating television state; and
  \item \texttt{tv.control} for session ownership and allowlisted TV commands.
\end{itemize}

Every request remains bound to the phone certificate's person and device.  Persona substitution,
missing scope, expired lease, revoked phone, unavailable adapter, locked shared screen or
unsupported command fails closed.

\subsection{Verification and Remaining Qualification}

Automated tests prove projection redaction, capability summarization, exclusive session
enforcement, nested possession signatures, locked-screen rejection, accepted TV control and the
distinction between infrared delivery and appliance-state verification.  Mobile-surface checks
cover separate Devices and TV experiences, observe-only discovery language, acquisition, remote
controls, application launch controls and the IoT verification warning.  At coded closure of
\texttt{MOB-0509}, all 465 AION Fabric tests passed and the optimized Next.js build generated
\texttt{/pilot/mobile} successfully.

Real-device qualification remains required after the Pilot runtime and phone lease are renewed:
the existing LG television must prove acquisition, observation, one reversible control and
automatic locking on the physical phone.  Additional television and appliance vendors remain
adapter- and hardware-dependent and are not claimed universal from the LG implementation alone.
No new device was enrolled and no live IoT command was sent during automated verification.
\endgroup
